\subsection*{Problem 1.3}

\subsubsection*{(a)}
If we \emph{consider} the \emph{suits} of the cards,
then there are $\binom{8}{2}\cdot\binom{6}{2}
\cdot\binom{4}{2} = 2520$ possible worlds.

\subsubsection*{(b)}
If we \emph{ignore} the \emph{suits} of the cards,
then there are $19$ possible worlds.
Consider the distributions of \emph{aces}
among the hands of each player and the desk:
\begin{gather*}
    (2,2,0,0) \rightarrow \binom{4}{2} = 6 \\
    (2,1,1,0) \rightarrow \binom{4}{1}\cdot\binom{3}{2} = 12 \\
    (1,1,1,1) \rightarrow 1
\end{gather*}
Hence the total is $6+12+1 = 19$ possible worlds.

\subsubsection*{(c)}
We can define the initial Kripke structure
(ignoring the \emph{suits}) as \krip{},
such that the states are the distributions of
\emph{aces} to Aneeket, Raj, me, and the desk
(respectively).
\begin{gather*}
    S = \{ a_1a_2a_3a_4 : 
        \forall i \in \{1,2,3,4\}:
            a_i \in \{0,1,2\}
            \land \sum_{i=1}^{4} a_i = 4
        \}, \\
    \text{Agents} = \{a,r,y\}, \\
    a_1a_2a_3a_4 \sim_{a} b_1b_2b_3b_4 
        \iff a_2 = b_2 \land a_3 = b_3, \\
    a_1a_2a_3a_4 \sim_{r} b_1b_2b_3b_4 
        \iff a_1 = b_1 \land a_3 = b_3, \\
    a_1a_2a_3a_4 \sim_{y} b_1b_2b_3b_4 
        \iff a_1 = b_1 \land a_2 = b_2,
    \tag{1.3(c1)}\label{eq:1.3(c1)} \\
    \forall n \in \{0,1,2\}: \\
    p_{a}^{(n)} = \{ n a_1 a_2 a_3 : 
        \forall i \in \{1,2,3\}: 
            a_i \in \{0,1,2\}
            \land \sum_{i=1}^{3} a_i = 4-n \}, \\
    p_{r}^{(n)} = \{ a_1 n a_2 a_3 : 
        \forall i \in \{1,2,3\}: 
            a_i \in \{0,1,2\}
            \land \sum_{i=1}^{3} a_i = 4-n \}, \\
    p_{y}^{(n)} = \{ a_1 a_2 n a_3 : 
        \forall i \in \{1,2,3\}: 
            a_i \in \{0,1,2\}
            \land \sum_{i=1}^{3} a_i = 4-n \}.
\end{gather*}

\subsubsection*{(d)}
I am holding \textbf{one ace and one eight}.

\subparagraph*{Intuitively:}
Since only two aces and two eights remain,
the possible cards in my hand are one ace and one eight,
two aces, or two eights. If I had two aces, then Raj
would know he must have two aces; similarly, if I had
two eights, then Aneeket would know he must have two aces.
Hence the only possible choice is one ace and one eight.

\subparagraph*{Considering the Kripke structure:}
We encode the announcements as
\[
\phi_i = \bigwedge_{n=0}^{2} \neg K_i p_{i}^{(n)}
\]
for agent $i$ being unable to determine his own cards.
The restricted Kripke structure after $\phi_a$ is
$M^{(1)} = M|_{\phi_a} = \langle S^{(1)}, \sim^{(1)}, 
V^{(1)} \rangle$,
where $S^{(1)} = [\![ \phi_a ]\!] = \{ s \in S : 
M,s \models \phi_a \}$.
\begin{align*}
    M,s \models \phi_a
    &\iff M,s \models \bigwedge_{n=0}^{2} 
        \neg K_a p_{a}^{(n)} \\
    &\iff \bigwedge_{n=0}^{2} M,s \models 
        \neg K_a p_{a}^{(n)} \\
    &\iff \bigwedge_{n=0}^{2} M,s \models 
        \neg \bigl( \forall t \in \sim_a(s) : 
            t \in p_{a}^{(n)} \bigr)
        \tag{1.3(c2)}\label{eq:1.3(c2)}
\end{align*}
Define
\[
a = \bigcup_{n \in \{0,2\}} \{ a_1 n n a_4 \in S \}
\]
representing the worlds where Raj and I 
both have either two aces or two eights.
From \eq{1.3(c2)} we observe that no world in 
$a$ satisfies the formula.
Intuitively, for any world in $a$,
 Aneeket would be able to know his cards, 
 so all those worlds should be eliminated after 
 his announcement.
Thus $S^{(1)} = [\![ \phi_a ]\!] = S \setminus a$, 
and the corresponding edges are removed from 
the Kripke structure.

Similarly, for Raj's announcement we define
\[
r = \bigcup_{n \in \{0,2\}} \{ n a_2 n a_4 \in S^{(0)} \}
\]
representing the worlds where Aneeket and I both 
have either two aces or two eights.
The formula \eq{1.3(c2)} is updated by replacing 
$\sim_a$ with $\sim_r^{(1)}$ and $M$ with $M^{(1)}$.
All these worlds are eliminated, giving $S^{(2)} = 
[\![ \phi_r ]\!] = S^{(1)} \setminus r$, 
with their corresponding edges removed.

\subsubsection*{(e)}
Again, I am holding \textbf{one ace and one eight}.

\subparagraph*{Intuitively:}
Since three aces and one eight remain,
the possible cards in my hand are one ace and 
one eight or two aces.
Initially I would not know my cards in any case.
If I had two aces, then Aneeket would not know his cards, 
but after Aneeket's announcement
Raj (who can see my and Aneeket's cards) 
would know that he must have one ace and one eight.
Because if Raj had two aces or two eights, then either 
I or Aneeket would have known our cards.
Hence the only possible hand is one ace and one eight.

\subparagraph*{Considering the Kripke structure:}
The encoding of the announcements is the same as in part~(d). 
Define
\[
y = \bigcup_{n \in \{0,2\}} \{ n n a_3 a_4 \in S \}
\]
representing the worlds where Aneeket and Raj both have 
either two aces or two eights.
The formula \eq{1.3(c2)} is updated by replacing $\sim_a$ 
with $\sim_y$.
With the same logic as above, for all worlds in $y$ 
I would be able to know my cards, so they are eliminated 
(along with their corresponding edges), yielding
$S^{(1)} = [\![ \phi_y ]\!] = S \setminus y$.
In the same way, we can define sets $a$ and $r$, 
and they will be eliminated after the announcements 
$\phi_a$ and $\phi_r$ respectively, with their corresponding edges removed.